\documentclass[11pt]{article}
\usepackage[T1]{fontenc}
\usepackage[utf8]{inputenc}
\usepackage{lmodern,microtype}
\usepackage[margin=1.05in]{geometry}
\usepackage{amsmath,amssymb,amsthm,mathtools}
\usepackage{booktabs,array}
\usepackage{xurl}
\usepackage[hidelinks]{hyperref}
\hypersetup{pdftitle={Clocks and First-Passage Measures in Tao's Collatz Theorem},pdfauthor={Kokuno Yumeto}}
\newtheorem{theorem}{Theorem}[section]
\newtheorem{proposition}[theorem]{Proposition}
\newtheorem{lemma}[theorem]{Lemma}
\newtheorem{corollary}[theorem]{Corollary}
\theoremstyle{definition}
\newtheorem{definition}[theorem]{Definition}
\theoremstyle{remark}
\newtheorem{remark}[theorem]{Remark}
\newcommand{\N}{\mathbb N}
\newcommand{\Z}{\mathbb Z}
\newcommand{\R}{\mathbb R}
\newcommand{\Opos}{\mathcal O}
\newcommand{\Prob}{\mathbb P}
\newcommand{\E}{\mathbb E}
\newcommand{\Col}{\mathrm C}
\newcommand{\Short}{\mathrm T}
\newcommand{\Syr}{\mathrm S}
\newcommand{\vnu}{\nu_2}
\newcommand{\dd}{\dagger}
\newcommand{\push}{_{*}}
\newcommand{\normone}[1]{\left\|#1\right\|_1}
\newcommand{\ind}{\mathbf 1}
\title{Clocks and First-Passage Measures\\in Tao's Collatz Theorem}
\author{Kokuno Yumeto}
\date{Preprint draft, 4 September 2026}
\begin{document}
\maketitle
\begin{abstract}
For every prescribed function tending to infinity, Tao proved that the
Collatz orbit minimum lies below that function on a set of starting
integers of logarithmic density one.
An exact treatment of the orbit clocks and sampling measures gives a
first-passage formulation of this result.  Valuation sums relate ordinary
Collatz time, division-counting time, and odd-return time.  With a separate
failure state, the first-passage pushforwards satisfy composition identities
that turn summable scale errors into a quantitative transport bound.  The
same estimate gives, along a common sequence of starting scales, limiting
first-entry laws at every threshold. These laws are exactly compatible
under further first passage, and their joint distributions at any finite
collection of thresholds converge with a common quantitative bound.
Applied to Tao's first-passage estimate, the transport bound
yields his fixed-threshold and diverging-threshold conclusions with exact
harmonic weights.  Two further estimates supply a multiplicative endpoint
buffer and a uniform coefficient bound in the first-passage argument.
An explicit reserve calculation verifies the margin used in the
$3$-adic separation estimate.
\end{abstract}

\section{Introduction}

Collatz descent is studied with three natural clocks: every application of
the usual map, every division by two, or every passage from one odd number
to the next.  These clocks give the same orbit minimum but different
stopping times.  Their exact relation becomes important when a descent
estimate is restricted to a finite time interval or when the distribution
of an orbit at its first visit to a smaller scale is used as the input to
another estimate.

The historical descent results of Terras and Everett, and the power-descent
results of Allouche and Korec, provide the setting for Tao's
almost-bounded theorem, as reviewed in his Section~1.1
\cite{Terras,Everett,Allouche,Korec,Tao7}.
Repeated descent requires control of more than the exceptional set at one
scale: the image of a typical set need not be typical for the measure used
at the next scale.  Tao addresses this difficulty by comparing first-passage
distributions.  Successive descent steps can then be combined through the
transport of measures, rather than through an independence assumption on
the successive values of a single orbit.

The clock identities in Section~\ref{sec:clocks} separate a change of time
parameter from a change of sampling measure.  Passing from arbitrary
starting integers to their odd parts, for example, requires the exact
dyadic fibre weights at each finite cutoff.  Section~\ref{sec:transport}
then treats first passage as a family of maps indexed by the destination
threshold.  A separate failure state makes their composition identities
valid on the entire domain.  The corresponding pushforwards contract
total variation, so discrepancies between adjacent scales can be added
before taking a limit.

The resulting transport estimate formulates the mechanism in Tao's
Section~3. Section~\ref{sec:consequences} constructs a family of limiting
first-entry laws at all thresholds from one common sequence of inputs.
The family satisfies exact first-passage compatibility. Its finite joint
laws obey the same total-variation error bound, independently of the number
of thresholds. The section also passes from a fixed-threshold estimate
to an arbitrary diverging threshold. The latter
tightness formulation is already present in Tao's Remark~1.4 and
Theorem~3.1.  Keeping the finite harmonic sums throughout identifies the
sampling law and the order of limits in each step.

Two estimates needed in the first-passage proof are developed in
Section~\ref{sec:repairs}.  A buffer in logarithmic time corresponds to a
multiplicative buffer in the initial value.  The coefficient estimate
requires a geometric weight that remains summable uniformly in the word
length.  Section~\ref{sec:versions} gives an explicit size bound for the
retained reserve in the $3$-adic separation argument and records the
finite-weight and threshold reductions to which it is adjacent.

The analytic input throughout is Tao's Proposition~1.11 in
arXiv:1909.03562v7, dated 16 July 2026 \cite{Tao7}, whose proof uses
Fourier--renewal estimates.  The deduction from that proposition and the
local estimates above are proved here.  The first arXiv version appeared
in September 2019, and the journal paper appeared in 2022
\cite{TaoJournal}. The exact compatible family is a distributional
consequence of the first-passage estimate; the quantitative orbit-minimum
bound remains Tao's Theorem~3.1.

\section{Three clocks and one orbit minimum}\label{sec:clocks}

Write $\N_0=\{0,1,\ldots\}$, $\N_+=\{1,2,\ldots\}$, and
$\Opos=\{1,3,5,\ldots\}$.  Throughout, logarithms are natural.
The three maps are
\[
 \Col(n)=\begin{cases}3n+1&n\text{ odd},\\ n/2&n\text{ even},\end{cases}
 \qquad
 \Short(n)=\begin{cases}(3n+1)/2&n\text{ odd},\\ n/2&n\text{ even},\end{cases}
\]
and $\Syr(m)=(3m+1)/2^{\vnu(3m+1)}$ for $m\in\Opos$.
The first two maps have domain and codomain $\N_+$, and the third has
domain and codomain $\Opos$.  The map $\Short$ counts one division by two
per step.  The map $\Syr$ counts one multiplication by three per step.
Tao makes this distinction in Section~1.2 of \cite{Tao7}; the shortened
clock also occurs explicitly in Terras and Korec \cite{Terras,Korec}.

For $N=2^a m$ with $m$ odd, define
\[
 m_j=\Syr^j(m),\qquad q_j=\vnu(3m_j+1),\qquad
 Q_0=0,\qquad Q_j=\sum_{i=0}^{j-1}q_i\quad(j\geq1).
\]
Every $q_j$ is a positive integer.  For a map $F$ on positive integers,
write $F_{\min}(n)=\min_{k\geq0}F^k(n)$; the minimum exists by
well-ordering, regardless of whether the orbit is bounded.

\begin{proposition}[Exact clock conversion]\label{prop:clocks}
For every $j\in\N_0$,
\begin{equation}\label{eq:clock}
 \Short^{a+Q_j}(N)=m_j=\Col^{a+j+Q_j}(N).
\end{equation}
The maps $j\mapsto a+Q_j$ and $j\mapsto a+j+Q_j$ are strictly
increasing bijections onto the respective sets of odd-visit times.
On these images the inverse maps assign to each odd-visit time its unique
index $j$ in the displayed equalities.
Furthermore,
\begin{equation}\label{eq:min}
 \Col_{\min}(2^a m)=\Short_{\min}(2^a m)=\Syr_{\min}(m).
\end{equation}
\end{proposition}
\begin{proof}
The first $a$ iterates under either full-domain map divide by two, reaching
$m$.  Starting at the full-clock time $a+j+Q_j$, the subsequent segment is
\[
 \Col^{a+j+Q_j+r}(N)=2^{q_j-r+1}m_{j+1}
 \quad(1\leq r\leq q_j+1).
\]
It has $q_j+1$ steps.  The shortened segment omits the first even value
$2^{q_j}m_{j+1}$ and has $q_j$ steps.  All intervening values are even,
and the terminal value is odd.  Induction proves \eqref{eq:clock} and
identifies exactly, with no additional visits, the odd-visit times.
Every deleted value is at least the following odd value, and the initial
segment is at least $m$.  Deletion therefore preserves the minimum.
The full sequence can be recovered from $a$ and the indexed odd orbit:
each $q_j$ is determined by $m_j$ and the displayed segments fill every
omitted time.  Repeated values on a periodic orbit do not affect this
indexed reconstruction.
\end{proof}

For example, the odd orbit beginning at $3$ is $3,5,1,1,\ldots$ with
initial valuations $1,4,2,\ldots$.  The shortened orbit visits these
states at times $0,1,5,7,\ldots$, whereas the full orbit visits them at
$0,2,7,10,\ldots$.  Thus a horizon of $j$ odd returns corresponds to
$Q_j$ shortened steps and $j+Q_j$ full steps.

The corresponding change of starting coordinates also changes weights.
Set
\[
 H(X)=\sum_{1\leq n\leq X}\frac1n,\qquad
 H_o(X)=\sum_{\substack{1\leq m\leq X\\m\text{ odd}}}\frac1m,
\]
with either sum zero when its range is empty.  Integral comparison gives
$H(X)=\log X+O(1)$ and $H_o(X)=\frac12\log X+O(1)$ as $X\to\infty$.

\begin{proposition}[Dyadic weight transport]\label{prop:dyadic}
The map $(a,m)\mapsto2^a m$ is a bijection
$\N_0\times\Opos\to\N_+$ with inverse
$n\mapsto(\vnu(n),n/2^{\vnu(n)})$.
For every $A\subseteq\Opos$, $a\in\N_0$, and $X\geq1$,
\begin{equation}\label{eq:weight}
 \sum_{\substack{n\leq X\\n\in2^aA}}\frac1n
 =2^{-a}\sum_{\substack{m\leq X/2^a\\m\in A}}\frac1m.
\end{equation}
Consequently the stratum $\vnu(n)=a$ has logarithmic density
$2^{-a-1}$, while
\begin{equation}\label{eq:dyadic-tail}
 \sum_{\substack{n\leq X\\\vnu(n)>A_0}}\frac1n
 =2^{-A_0-1}H(X/2^{A_0+1})\qquad(A_0\in\N_0).
\end{equation}
\end{proposition}
\begin{proof}
Unique factorization by the largest power of two gives the bijection,
whose fibres are singletons.  Substitution in the finite sum proves
\eqref{eq:weight}.  For fixed $a$ its right-hand side with $A=\Opos$ is
$2^{-a-1}\log X+O_a(1)$.  Division by $H(X)$ proves the density claim.
The tail consists exactly of multiples of $2^{A_0+1}$; substituting
$n=2^{A_0+1}r$ proves \eqref{eq:dyadic-tail}.
\end{proof}

The odd-part projection $\pi(n)=n/2^{\vnu(n)}$ has fibre
$\{2^a m:a\geq0\}$ over $m$.  Retaining $a$ recovers the starting
integer; \eqref{eq:clock} supplies the corresponding time change.
This time change is essential: for example, $\pi(\Col(6))=3$, whereas
$\Syr(\pi(6))=5$.

There is also a useful comparison at a fixed cutoff.  For an integer
$X\geq1$, let $\rho_X(n)=1/(nH(X))$ on $1\leq n\leq X$ and let
$\rho_X^o(m)=1/(mH_o(X))$ on odd $m\leq X$.

\begin{proposition}[Finite-cutoff sampling comparison]\label{prop:sampling}
The odd-part pushforward obeys
\[
 \sup_{E\subseteq\Opos}|(\pi_*\rho_X)(E)-\rho_X^o(E)|
 \leq\frac{H(X)-H(\lfloor X/2\rfloor)}{H(X)}.
\]
In particular its $\ell^1$ distance from $\rho_X^o$ is $O(1/\log X)$
as $X\to\infty$.  Upper and lower logarithmic densities of
$\pi^{-1}(E)$ therefore equal the respective relative odd logarithmic
densities of $E$, whether or not either density exists as a limit.
\end{proposition}
\begin{proof}
For odd $m\leq X$, let $r_m$ be the unique integer with
$2^{r_m}m\leq X<2^{r_m+1}m$.  The fibre mass is exactly
$(2-2^{-r_m})/(mH(X))$.  Put
\[
 B_E=\sum_{\substack{m\leq X\\m\in E}}\frac1m,
 \qquad D_E=\sum_{\substack{m\leq X\\m\in E}}\frac{2^{-r_m}}m,
 \qquad D=D_{\Opos}.
\]
The fibre partition gives $H(X)=2H_o(X)-D$; separating the even terms
also gives $D=H(X)-H(\lfloor X/2\rfloor)$, with $H(0)=0$.
Writing $q=B_E/H_o(X)$, direct subtraction yields
\[
 (\pi_*\rho_X)(E)-\rho_X^o(E)=\frac{qD-D_E}{H(X)}.
\]
Since $0\leq q\leq1$ and $0\leq D_E\leq D$, this proves the event
bound.  The $\ell^1$ distance of probability measures is twice the
supremum event difference: the positive and negative parts of their
difference have equal mass.  Finally $D\to\log2$ and $H(X)\sim\log X$.
The assertions about upper and lower densities follow by taking limsups
and liminfs in the finite bound.
\end{proof}

\section{First passage as an exact transport map}\label{sec:transport}

Let $D$ be a nonempty subset of $\N_+$ and $F:D\to D$ any map.
For a real $x\geq1$, define the first entrance time into $D\cap[1,x]$ by
\[
 t_x(n)=\inf\{j\in\N_0:F^j(n)\leq x\},\qquad\inf\varnothing=+\infty.
\]
Adjoin a point $\dd\notin D$, and define $p_x:D\cup\{\dd\}\to
D\cup\{\dd\}$ by
\[
 p_x(n)=\begin{cases}F^{t_x(n)}(n)&t_x(n)<\infty,\\
 \dd&t_x(n)=\infty,\end{cases}\qquad p_x(\dd)=\dd.
\]
The absorbing point $\dd$ records trajectories that never enter
$D\cap[1,x]$.  The image of $p_x$ is contained in
$(D\cap[1,x])\cup\{\dd\}$, and every point of this set is fixed by $p_x$.
The fibre $p_x^{-1}(\dd)\cap D$ consists exactly of their starting values
and may be empty.

\begin{lemma}[Nested passage]\label{lem:nested}
For $1\leq x\leq y$,
\begin{equation}\label{eq:nested}
 p_x\circ p_y=p_x=p_y\circ p_x.
\end{equation}
When $t_x(n)<\infty$, both hitting times are finite and
\[
 t_x(n)=t_y(n)+t_x(p_y(n)),\qquad
 F^{k}(p_x(n))=F^{k+t_x(p_y(n))}(p_y(n))\quad(k\geq0).
\]
\end{lemma}
\begin{proof}
If the orbit never reaches $y$, it cannot reach $x$, and the first
identity gives $\dd$ on both sides.  If it reaches $y$, the orbit prior
to that first visit is above $y$, hence above $x$.  Searching for the
first visit to $x$ can therefore begin at $p_y(n)$, giving the hitting-time
identity when finite and $\dd$ on both sides otherwise.  The equality
$p_y p_x=p_x$ follows because a successful output of $p_x$ is at most
$x\leq y$, and $\dd$ is fixed.  The identity of indexed tails follows
from the semigroup law for $F$.
\end{proof}

For a signed summable measure $\lambda$ on $D\cup\{\dd\}$, define its
pushforward $K_x\lambda=(p_x)_*\lambda$, explicitly
\[
 (K_x\lambda)(z)=\sum_{n:p_x(n)=z}\lambda(n).
\]
These maps preserve total mass, preserve nonnegativity, and contract
$\ell^1$: $\normone{K_x\lambda}\leq\normone{\lambda}$.  Indeed, take
absolute values inside each fibre sum and sum over the disjoint fibres.
Lemma~\ref{lem:nested} gives
\begin{equation}\label{eq:kernel}
 K_xK_y=K_x=K_yK_x\quad(x\leq y),\qquad K_x^2=K_x.
\end{equation}
The kernel of $K_x$ consists exactly of signed measures whose sum on
each fibre of $p_x$ is zero.  In particular,
$K_x(\delta_n-\delta_{n'})=0$ whenever $p_x(n)=p_x(n')$.

\begin{proposition}[Finite-chain transport bound]\label{prop:chain}
Let $1\leq x_0\leq\cdots\leq x_J$ and let $\nu_0,\ldots,\nu_J$
be probability measures satisfying $K_{x_i}\nu_i=\nu_i$.
Put $e_i=\normone{\nu_i-K_{x_i}\nu_{i+1}}$ for $0\leq i<J$.
Then
\begin{equation}\label{eq:chain}
 \normone{\nu_0-K_{x_0}\nu_J}\leq\sum_{i=0}^{J-1}e_i.
\end{equation}
\end{proposition}
\begin{proof}
For $J=0$ both sides are zero.  For $J\geq1$, the exact telescoping
identity is
\[
 \nu_0-K_{x_0}\nu_J=
 \sum_{i=0}^{J-1}K_{x_0}(\nu_i-K_{x_i}\nu_{i+1}).
\]
Here $K_{x_0}\nu_0=\nu_0$ and $K_{x_0}K_{x_i}=K_{x_0}$ by
\eqref{eq:kernel}.  The triangle inequality and contraction prove
\eqref{eq:chain}.
\end{proof}

The following scale formulation records exactly the summability used in
Tao's reduction.  Fix $\alpha>1$.  Suppose $\mu_s$, for $s\geq1$, is a probability
measure on $D\cap[s,s^\alpha]$ whenever that set is nonempty.  Choose
$s_*\geq2$ so every required block is nonempty for $s\geq s_*$.
For such $s$, set
\[
 \nu_s=K_s\mu_{s^\alpha},\qquad d(s)=\nu_s(\{\dd\}),\qquad
 e(s)=\normone{\nu_s-K_s\nu_{s^\alpha}}.
\]
For $R\geq s_*$ define the possibly infinite quantity
\begin{equation}\label{eq:error-tail}
 \mathcal E(R)=\sup_{u\geq R}
 \left(d(u)+\sum_{i=0}^{\infty}e(u^{\alpha^i})\right).
\end{equation}
Here $d(s)$ is the failure mass at scale $s$, while $e(s)$ measures the
change in the entrance law between two successive starting scales.

\begin{theorem}[Summable transport implies tight orbit minima]
\label{thm:transport}
For $K\geq s_*^{\alpha^3}$ and any nonempty block supporting $\mu_s$,
\begin{equation}\label{eq:abstract-tight}
 \mu_s\{n:F_{\min}(n)>K\}\leq\mathcal E(K^{1/\alpha^3}).
\end{equation}
In particular, if $\mathcal E(R)\to0$ as $R\to\infty$, the laws of
$F_{\min}$ under these block measures are uniformly tight.
\end{theorem}
\begin{proof}
Let $A_K=\{n\in D:F_{\min}(n)\leq K\}$, so $\dd\notin A_K$.
A successful passage landing in $A_K$ certifies that the original
orbit enters $A_K$.  Thus $p_u^{-1}(A_K)\subseteq A_K$ on the extended
space, for every $u$.

If $s^\alpha\leq K$, the left side of \eqref{eq:abstract-tight} is zero.
Otherwise choose $J\geq1$ minimal such that
$y=s^{\alpha^{-J}}<K^{1/\alpha}$.  Then
\[
 K^{1/\alpha^2}\leq y<K^{1/\alpha},\qquad
 u=y^{1/\alpha}\geq K^{1/\alpha^3}\geq s_*,\qquad
 u^{\alpha^J}=s^{1/\alpha}.
\]
The successful part of $\nu_u$ is supported below $u<K$, so
$\nu_u(A_K)=1-d(u)$.
For each intermediate scale $v=u^{\alpha^i}$,
\[
 \nu_{v^\alpha}(A_K)
 \geq (K_v\nu_{v^\alpha})(A_K)
 \geq \nu_v(A_K)-e(v).
\]
Summing for $0\leq i<J$ gives
\[
 \nu_{s^{1/\alpha}}(A_K)
 \geq1-d(u)-\sum_{i=0}^{J-1}e(u^{\alpha^i}).
\]
But $\nu_{s^{1/\alpha}}=K_{s^{1/\alpha}}\mu_s$, and
$p_{s^{1/\alpha}}^{-1}(A_K)\subseteq A_K$, so
$\mu_s(A_K)$ is at least this large.  Formula \eqref{eq:error-tail}
proves the bound.
\end{proof}

For example, if $d(s)\leq Ds^{-b}$ and $e(s)\leq E(\log s)^{-c}$,
with $D,E,b,c>0$, then
\begin{equation}\label{eq:power-error}
 \mathcal E(R)\leq DR^{-b}
 +\frac{E}{1-\alpha^{-c}}(\log R)^{-c}.
\end{equation}
This follows by summing the geometric series $\sum_{i\geq0}\alpha^{-ci}$.
The same theorem also accepts, for instance,
$e(s)=O((\log\log s)^{-1-c})$ for large $s$: for sufficiently large
$R>e$, writing
$h=\log\alpha>0$ and $v=\log\log R$, the sum is bounded by
$v^{-1-c}+v^{-c}/(ch)$ by integral comparison.  Tao already notes this
weaker sufficient error regime after Proposition~1.11 \cite{Tao7}.

\section{Application to Tao's theorem}\label{sec:consequences}

Take $D=\Opos$, $F=\Syr$, and $\alpha=1001/1000$.  Let
\[
 R_s=\Opos\cap[s,s^\alpha],\qquad
 h_s=\sum_{n\in R_s}\frac1n,\qquad
 \mu_s(\{n\})=\frac{1}{h_s n}\quad(n\in R_s),
\]
whenever $R_s\ne\varnothing$.  The measure is zero outside $R_s$.
Choose $s_*$ sufficiently large that all blocks used below are nonempty.

\begin{theorem}[Tao, Proposition 1.11]\label{thm:tao-input}
There are absolute constants $B,c>0$ such that, for all sufficiently
large $x$ and $y\in\{x^\alpha,x^{\alpha^2}\}$, an input $\mathbf N_y$
with law $\mu_y$ satisfies
\[
 \Prob(t_x(\mathbf N_y)=\infty)\leq Bx^{-c}.
\]
Let $\widehat p_x(n)$ agree with $p_x(n)$ on successful passage and
equal $1$ on failed passage, as in Tao's convention.  Then
\begin{equation}\label{eq:tao-tv}
 \normone{(\widehat p_x)_*\mu_{x^\alpha}
 -(\widehat p_x)_*\mu_{x^{\alpha^2}}}
 \leq B(\log x)^{-c}.
\end{equation}
\end{theorem}

The norm in \eqref{eq:tao-tv} is the unhalved $\ell^1$ norm used in
\cite{Tao7}.  Tao proves this proposition using the valuation law and
fine-scale mixing of Syracuse random variables, obtained through Fourier
estimates and a two-dimensional renewal argument.  The estimates in
Section~\ref{sec:repairs} supply the endpoint margin and uniform
coefficient bound used in his Section~5.

The change to $\dd$ costs only the known failure probability.  Indeed,
if $q$ merges $\dd$ with $1$ and fixes the other points, then for two
probability measures $\lambda,\lambda'$ on $\Opos\cup\{\dd\}$,
\begin{equation}\label{eq:unmerge}
 \normone{\lambda-\lambda'}
 \leq\normone{q_*\lambda-q_*\lambda'}
 +2|\lambda(\dd)-\lambda'(\dd)|.
\end{equation}
To check it, only the coordinates $1$ and $\dd$ need attention:
$|a|+|b|\leq|a+b|+2|b|$ for their respective differences.
By \eqref{eq:kernel},
$K_x\nu_{x^\alpha}=K_x\mu_{x^{\alpha^2}}$.
Thus Theorem~\ref{thm:tao-input} proves, after increasing a constant,
\begin{equation}\label{eq:input-completed}
 d(s)\leq Bs^{-c},\qquad e(s)\leq B'(\log s)^{-c}
 \quad(s\geq s_*).
\end{equation}
These measures retain their failure mass and the original denominators $h_s$.

\subsection{An exactly compatible family of first-entry laws}

Fix one base $b\geq s_*$ and put $t_j=b^{\alpha^j}$ for $j\geq0$.
For every real threshold $x\geq1$, write
\[
 E_x=(\Opos\cap[1,x])\cup\{\dd\},\qquad
 L_c=\frac{B'}{1-\alpha^{-c}},
\]
and define the probability measure
\begin{equation}\label{eq:lattice-law}
 \lambda_{x,j}^{(b)}=K_x\mu_{t_{j+1}}\quad\hbox{on }E_x.
\end{equation}
The input sequence is the same for every threshold. These laws describe
the first entry into $[1,x]$, including failure, rather than the position
after a prescribed number of iterations.

\begin{corollary}[Coherent first-entry limits]\label{cor:limit}
For each $x\geq1$, the measures $\lambda_{x,j}^{(b)}$ converge in total
variation to a probability measure $\lambda_x^{\infty,(b)}$ on $E_x$.
For every $j\geq0$ and $1\leq x\leq t_j$,
\begin{equation}\label{eq:lattice-rate}
 \normone{\lambda_{x,j}^{(b)}-\lambda_x^{\infty,(b)}}
 \leq L_c(\log t_j)^{-c}.
\end{equation}
The limiting family is exactly compatible at every pair of thresholds:
\begin{equation}\label{eq:lattice-compatible}
 K_x\lambda_y^{\infty,(b)}=\lambda_x^{\infty,(b)}
 \qquad(1\leq x\leq y).
\end{equation}
For $x\geq b$, its failure mass obeys
\begin{equation}\label{eq:all-threshold-failure}
 \lambda_x^{\infty,(b)}(\dd)
 \leq Bx^{-c/\alpha}+L_c\alpha^c(\log x)^{-c}.
\end{equation}
\end{corollary}

\begin{proof}
When $x\leq t_j$, nested passage gives
\[
 \lambda_{x,j}^{(b)}=K_x\nu_{t_j},\qquad
 \lambda_{x,j+1}^{(b)}=K_xK_{t_j}\nu_{t_{j+1}}.
\]
Hence contraction and \eqref{eq:input-completed} imply
\[
 \normone{\lambda_{x,j+1}^{(b)}-\lambda_{x,j}^{(b)}}
 \leq e(t_j)\leq B'(\alpha^j\log b)^{-c}.
\]
For each fixed $x$, these inequalities apply from some index onward.
The geometric tail makes the sequence Cauchy in the finite-dimensional
space $\ell^1(E_x)$. Its limit is nonnegative with total mass one.
Summing the tail from $j$ proves \eqref{eq:lattice-rate}, uniformly for
all $x\leq t_j$.

For $x\leq y$, the identity
$K_x\lambda_{y,j}^{(b)}=\lambda_{x,j}^{(b)}$ holds at every finite $j$.
Taking limits under the contraction $K_x$ proves
\eqref{eq:lattice-compatible}. In particular, the thresholds need not
belong to the sequence $(t_j)$.

For $x\geq b$, choose the unique $k\geq0$ with
$r=t_k\leq x<t_{k+1}=r^\alpha$. Then $r>x^{1/\alpha}$, and
$\lambda_{r,k}^{(b)}=\nu_r$. Equations \eqref{eq:input-completed} and
\eqref{eq:lattice-rate} give
\[
 \lambda_r^{\infty,(b)}(\dd)\leq Br^{-c}+L_c(\log r)^{-c}.
\]
Since $p_r(\dd)=\dd$, compatibility gives
$\lambda_x^{\infty,(b)}(\dd)\leq\lambda_r^{\infty,(b)}(\dd)$.
Substitution of $r>x^{1/\alpha}$ proves
\eqref{eq:all-threshold-failure}.
\end{proof}

The base $b$ records the logarithmic scale phase. Replacing it by
$b^{\alpha^k}$, $k\geq0$, deletes finitely many initial inputs and
therefore gives exactly the same limiting family. Equality for arbitrary
different bases is not supplied by this argument. For a threshold
$x=t_k$, the limit also equals the threshold-based limit obtained from
inputs $\mu_{x^{\alpha^{j+1}}}$, since that sequence is a tail of the
common input sequence.

\begin{corollary}[Joint first-entry distributions]\label{cor:joint-entry}
Fix $1\leq x_1\leq\cdots\leq x_r$, with $r\geq1$, and let
$\mathbf N_j$ have law $\mu_{t_{j+1}}$. The joint law of
\[
 (p_{x_1}(\mathbf N_j),\ldots,p_{x_r}(\mathbf N_j))
\]
converges to a probability measure $\Lambda_{x_1,\ldots,x_r}^{(b)}$.
It is supported on the finite set
\[
 \mathcal C=\{(z_1,\ldots,z_r)\in\textstyle\prod_i E_{x_i}:
                p_{x_i}(z_{i+1})=z_i\ (1\leq i<r)\}.
\]
For every $j$ with $t_j\geq x_r$, its unhalved total-variation error
is at most $L_c(\log t_j)^{-c}$, with no factor depending on $r$.
The marginal at $x_i$ is $\lambda_{x_i}^{\infty,(b)}$.
\end{corollary}

\begin{proof}
Define
\[
 J:E_{x_r}\longrightarrow\mathcal C,\qquad
 J(z)=(p_{x_1}(z),\ldots,p_{x_r}(z)).
\]
Nested passage proves that $J$ takes values in $\mathcal C$. Its last
coordinate is $z$, so it is injective. Conversely, repeated use of the
defining equations of $\mathcal C$ gives $z_i=p_{x_i}(z_r)$; thus the
inverse of $J$ is projection to the last coordinate. In particular $J$
has singleton fibres and loses no information. Its pushforward preserves
the $\ell^1$ norm of signed measures, because distinct source atoms have
distinct images.

The finite joint law is exactly $J_*\lambda_{x_r,j}^{(b)}$.
Set $\Lambda_{x_1,\ldots,x_r}^{(b)}=J_*\lambda_{x_r}^{\infty,(b)}$.
The isometry and \eqref{eq:lattice-rate} give the stated error bound.
Projection onto coordinate $i$ and
\eqref{eq:lattice-compatible} identify the marginal.
\end{proof}

Tao's Proposition~1.11 compares two successive starting scales. The
corollaries above turn its summable error into exact compatible
distributions, including simultaneous first-entry observations. The
compatibility maps are the first-passage maps $K_x$; no invariance under
one application of $\Syr$ is required.

\subsection{The quantitative fixed-threshold theorem}

\begin{theorem}[Tao's fixed-threshold theorem]
\label{thm:fixed}
There are absolute constants $C,c>0$ such that, for all real $K,X\geq2$,
\begin{align}
 \sum_{\substack{m\leq X,\ m\in\Opos\\\Syr_{\min}(m)>K}}\frac1m
 &\leq C\frac{\log X}{(\log K)^c},\label{eq:odd-fixed}\\
 \sum_{\substack{n\leq X\\\Col_{\min}(n)>K}}\frac1n
 &\leq C\frac{\log X}{(\log K)^c}.\label{eq:full-fixed}
\end{align}
\end{theorem}
\begin{proof}
Equations \eqref{eq:input-completed} and \eqref{eq:power-error}, followed
by Theorem~\ref{thm:transport}, give
\[
 \mu_s(\Syr_{\min}>K)\leq C_1(\log K)^{-c}
\]
for all sufficiently large $K$, uniformly over nonempty blocks.  Bounded
$K\geq2$ is included by increasing $C_1$.
For $s\geq2$, integral comparison gives $h_s\leq C_\alpha\log s$.
If $1<s<2$, the block contains no odd integer, since $s^\alpha<2^\alpha<3$.
Hence every relevant block obeys
\begin{equation}\label{eq:block-harmonic}
 \sum_{\substack{m\in R_s\\\Syr_{\min}(m)>K}}\frac1m
 \leq C_2(\log s)(\log K)^{-c},
\end{equation}
with empty blocks contributing zero.

If $X\leq K$, the sum in \eqref{eq:odd-fixed} is zero, since the orbit
contains its starting value.  Otherwise let $z_i=X^{\alpha^{-i}}$, and
let $J\geq1$ be minimal with $z_J\leq K$.  The blocks
$[z_i,z_i^\alpha]$, $1\leq i\leq J$, cover $[z_J,X]$.
Every bad starting value exceeds $K$, so the remaining prefix contributes
nothing.  Possible duplicated endpoints only increase an upper bound.
Since
\[
 \sum_{i=1}^J\log z_i
 =\log X\sum_{i=1}^J\alpha^{-i}
 \leq\frac{\log X}{\alpha-1},
\]
summing \eqref{eq:block-harmonic} proves \eqref{eq:odd-fixed}.

For the full-domain map, Propositions~\ref{prop:clocks} and
\ref{prop:dyadic} give the finite identity
\[
 \sum_{\substack{n\leq X\\\Col_{\min}(n)>K}}\frac1n
 =\sum_{a\geq0}2^{-a}
   \sum_{\substack{m\leq X/2^a,\ m\in\Opos\\\Syr_{\min}(m)>K}}\frac1m.
\]
Only finitely many terms are nonzero.  If $X/2^a<2$, the inner bad sum is
zero.  For the others, apply \eqref{eq:odd-fixed} and use
\[
 \sum_{a:X/2^a\geq2}2^{-a}\log(X/2^a)
 \leq\log X\sum_{a\geq0}2^{-a}=2\log X.
\]
This proves \eqref{eq:full-fixed}.
\end{proof}

\subsection{Diverging thresholds and the exact meaning of tightness}

\begin{corollary}[Tao's almost-bounded theorem]\label{cor:tao}
For every $f:\N_+\to\R$ with $f(n)\to+\infty$,
\[
 \lim_{X\to\infty}\frac1{H(X)}
 \sum_{\substack{n\leq X\\\Col_{\min}(n)\geq f(n)}}\frac1n=0.
\]
Thus $\Col_{\min}(n)<f(n)$ on a set of logarithmic density one.
The analogous statement on $\Opos$, with $H_o(X)$ and $\Syr$, also holds.
\end{corollary}
\begin{proof}
Fix an integer $K\geq2$.  There is $Y_K$ such that $f(n)>K$ for
$n\geq Y_K$.  Therefore
\[
 \{n:\Col_{\min}(n)\geq f(n)\}
 \subseteq[1,Y_K]\cup\{n:\Col_{\min}(n)>K\}.
\]
Theorem~\ref{thm:fixed} bounds its harmonic mass up to $X$ by
$H(Y_K)+C\log X/(\log K)^c$.  Divide by $H(X)$ and first let
$X\to\infty$ at fixed $K$.  The resulting limsup is at most
$C/(\log K)^c$, which tends to zero as $K\to\infty$.
The odd case is identical, using $H_o(X)\sim\frac12\log X$.
The argument allows nonmonotone $f$ with arbitrarily slow divergence.
Its negative values occur only in a finite prefix, and the equality part
of the bad event is included throughout.
\end{proof}

The finite-prefix argument gives a general tightness criterion for
weighted counting measures.

\begin{proposition}[Weighted tightness equivalence]\label{prop:tightness}
Let $D\subseteq\N_+$ be infinite, let $w(n)>0$, and suppose
$W(X)=\sum_{n\in D,n\leq X}w(n)\to\infty$.
Let $g:D\to[0,\infty)$ be finite-valued and let $\rho_X$ be the
probability proportional to $w$ on the nonempty prefix $D\cap[1,X]$.
The following assertions are equivalent:
\begin{enumerate}
\item For every real-valued $f(n)\to\infty$, $\rho_X(g\geq f)\to0$.
\item $\displaystyle\lim_{K\to\infty}\limsup_{X\to\infty}\rho_X(g>K)=0$.
\item $\displaystyle\lim_{K\to\infty}\sup_X\rho_X(g>K)=0$.
\end{enumerate}
The supremum is over nonempty prefixes.
\end{proposition}
\begin{proof}
The finite-prefix argument in Corollary~\ref{cor:tao}, with $W$ in place
of $H$, proves (2)$\Rightarrow$(1).
If (2) fails, choose $\varepsilon>0$ such that for every integer $j\geq1$
there are arbitrarily large $X$ with $\rho_X(g>j)\geq\varepsilon$.
Start with an integer $X_0\geq\min D$.
Choose integers $X_j>X_{j-1}$ recursively with this property and with
$W(X_{j-1})/W(X_j)\leq\varepsilon/2$.  Set $f(n)=j$ for
$X_{j-1}<n\leq X_j$, assigning any finite value on the initial prefix.
This tends to infinity.  At $X_j$, removal of the earlier prefix leaves
mass at least $\varepsilon/2$ on which $g>j=f$, contradicting (1).

To prove (2)$\Rightarrow$(3), choose $K_0$ so the limsup in (2) is less
than $\varepsilon/2$.  For some $X_0$, all $X\geq X_0$ have
$\rho_X(g>K_0)\leq\varepsilon$.  Increase $K_0$ to exceed every value
of $g$ on the finite prefix up to $X_0$.  All earlier tails are now empty,
and the later bound persists.  Finally (3)$\Rightarrow$(2) is immediate.
\end{proof}

For $w(n)=1/n$ and $g=\Col_{\min}$ or $\Syr_{\min}$, this gives the
qualitative equivalence stated by Tao.  The quantitative rate in
Theorem~\ref{thm:fixed} uses his first-passage estimate in addition to
this equivalence.

\section{Endpoint margins and uniform coefficient estimates}\label{sec:repairs}

Two estimates in the proof of Tao's Proposition~5.2 control the location
of the first-passage time and the effect of fine-scale mixing on the
first-entry distribution.  Retain his parameters
\begin{equation}\label{eq:local-parameters}
 L=\log x,\quad \alpha=1.001,\quad \beta=\log(4/3),\quad
 n_0=\left\lfloor\frac{L}{10\log2}\right\rfloor,\quad
 m_0=\left\lfloor\frac{L}{100000}\right\rfloor.
\end{equation}
Here $x$ is sufficiently large; in particular $m_0\geq1$.

\subsection{A buffer in time is multiplicative in the initial value}

In the proof of Proposition~5.2, Tao obtains an event $G_x$ with
$\Prob(G_x^c)=O(L^{-10})$ on which, for
$\mathbf N_y\sim\mu_y$ and $y\in\{x^\alpha,x^{\alpha^2}\}$,
\begin{equation}\label{eq:orbit-log}
 \left|\log\frac{\Syr^j(\mathbf N_y)}{\mathbf N_y}
          +j\beta\right|\leq C_0L^{3/5}
 \quad(0\leq j\leq n_0).
\end{equation}
The desired time interval is
\[
 I_y=\left[\frac{\log(y/x)}\beta+L^{4/5},\quad
             \frac{\log(y^\alpha/x)}\beta-L^{4/5}\right].
\]
The input restriction printed at this point in v7 adds and subtracts
$2L^{4/5}$ at the endpoints $y,y^\alpha$.
An additive displacement of that size produces logarithmic displacement
$O(L^{4/5}/y)$ at the lower endpoint, not the required time margin.
A multiplicative restriction supplies the required margin.

\begin{proposition}[Multiplicative endpoint estimate]\label{prop:buffer}
Set
\[
 J_y=[y\exp(L^{4/5}),\ y^\alpha\exp(-L^{4/5})].
\]
For all sufficiently large $x$ this is nonempty,
\[
 G_x\cap\{\mathbf N_y\in J_y\}\subseteq\{t_x(\mathbf N_y)\in I_y\},
 \qquad
 \Prob(t_x(\mathbf N_y)\in I_y)=1-O(L^{-1/5}),
\]
uniformly for the two indicated values of $y$.
\end{proposition}
\begin{proof}
Since $\mathbf N_y\leq x^{\alpha^3}$, \eqref{eq:orbit-log} gives
\[
 \log\frac{\Syr^{n_0}(\mathbf N_y)}x
 \leq(\alpha^3-1)L-\beta n_0+C_0L^{3/5}<0
\]
for large $x$.  For example $\alpha^3-1<1/300$ and
$\beta/(10\log2)>1/40$, whereas the floor costs only a constant.
Since $\mathbf N_y>x$, the first crossing is therefore at an integer
$1\leq t\leq n_0$.  Applying \eqref{eq:orbit-log} at $t$ and $t-1$
gives
\begin{equation}\label{eq:time-log}
 \left|t-\frac{\log(\mathbf N_y/x)}\beta\right|
 \leq C_1L^{3/5}.
\end{equation}
Indeed the inequalities $\Syr^t(\mathbf N_y)\leq x$
and $\Syr^{t-1}(\mathbf N_y)>x$ bound $t$ from below and above;
the additional one-step term is absorbed into $C_1$.

If $\mathbf N_y\in J_y$, the central expression in
\eqref{eq:time-log} lies inside $I_y$ with extra margin
$(\beta^{-1}-1)L^{4/5}$ on both sides.  Since $0<\beta<1$, this
dominates $C_1L^{3/5}$.  This proves the event inclusion.

For $1\leq a<b$, integral comparison along the odd progression gives
\[
 \sum_{\substack{a\leq m\leq b\\m\text{ odd}}}\frac1m
 =\frac12\log(b/a)+O(1/a),
\]
uniformly in the real endpoints.  Thus the block denominator is
$\frac12(\alpha-1)\log y+O(1/y)$, whereas the two excluded pieces have
total harmonic mass $L^{4/5}+O(1/y)$.  Their probability is
$O(L^{-1/5})$.  The logarithmic width of $J_y$ is
$(\alpha-1)\log y-2L^{4/5}>0$.  Finally add
$\Prob(G_x^c)=O(L^{-10})$.
\end{proof}

Thus the multiplicative buffer gives the same interval $I_y$ and
exceptional probability $O(L^{-1/5})$ required in the subsequent argument.

\subsection{A uniform coefficient bound over valuation words}

For a positive tuple $\mathbf a=(a_1,\ldots,a_m)$ put
\[
 A=a_1+\cdots+a_m,\qquad b=a_m,\qquad
 F_m(\mathbf a)=\sum_{i=1}^m3^{m-i}2^{-(a_i+\cdots+a_m)},
 \quad F_0=0.
\]
Composition of the maps $z\mapsto(3z+1)/2^{a_i}$ gives
\begin{equation}\label{eq:affine}
 \Syr^m(M)=3^m2^{-A}M+F_m(\mathbf a)
\end{equation}
when $\mathbf a$ is the actual valuation tuple of $M$.
Induction proves the formula: the next map multiplies every existing term
by $3/2^{a_{m+1}}$ and adds $2^{-a_{m+1}}$.
In particular $0\leq F_m\leq3^m$, a bound uniform in the tuple.
For the latter bound, replace every negative power of two by at most
$1/2$ and sum the finite geometric series.

For an arbitrary $E\subseteq\Opos\cap[1,x]$ set
\[
 E'=\left\{M\in\Opos:\begin{array}{l}
 t_x(M)=m_0,\quad p_x(M)\in E,\\
 e^{-L^{7/10}}(4/3)^{m_0}x\leq M
       \leq e^{L^{7/10}}(4/3)^{m_0}x
 \end{array}\right\}.
\]
This is the same set as in the approximate formula of
\cite[Proposition 5.2]{Tao7}; successful passage makes its meaning
independent of the failure convention.  For $n\in I_y\cap\Z$ write
$m=m_0$, $r=n-m\geq0$, and define on $\Z/3^r\Z$
\begin{equation}\label{eq:cn}
 c_n(z)=3^r\sum_{\substack{M\in E'\\M\equiv z\pmod{3^r}}}\frac1M.
\end{equation}
The source parameters ensure $m\leq n\leq n_0$ for sufficiently large
$x$.  The case $r=0$ is allowed in the estimate below, with modulus one.

\begin{proposition}[Uniform coefficient bound]\label{prop:cn}
There is an absolute constant $C$ such that $0\leq c_n(z)\leq C$
for all sufficiently large $x$, both choices of $y$, all
$n\in I_y\cap\Z$, all residue classes $z$, and every set $E$ above.
\end{proposition}
\begin{proof}
Partition \eqref{eq:cn} by the unique length-$m$ valuation tuple of $M$,
and call a tuple's contribution $c_{n,\mathbf a}(z)$.  At its first
crossing, the exact affine formulas give
\[
 3^m2^{-A}M+F_m\leq x
 <3^{m-1}2^{-(A-b)}M+F_{m-1}.
\]
The coefficient $3^{m-1}$ corresponds to the preceding $m-1$ iterates.
Uniformly $F_{m-1}\leq3^{m-1}=o(x)$, so every such $M$ lies between
\begin{equation}\label{eq:M-bounds}
 M_0=3^{-m}2^{A-b}x\quad\hbox{and}\quad M_1=3^{-m}2^A x
\end{equation}
for large $x$.  Indeed the upper bound follows by discarding $F_m\geq0$,
and the lower follows from $x-F_{m-1}>x/3$.

The integer $B_{\mathbf a}=2^A F_m$ is odd: its first term is
$3^{m-1}$ and all its later terms are even.  The terminal value in
\eqref{eq:affine} is odd, hence
\[
 3^mM+B_{\mathbf a}\equiv2^A\pmod{2^{A+1}}.
\]
Since $3^m$ is a unit modulo $2^{A+1}$, this requires a single class of
$M$ modulo $2^{A+1}$.  Combining it with $M\equiv z\pmod{3^r}$ gives
a single class modulo $q=2^{A+1}3^r$ by the Chinese remainder theorem.
Only this necessary restriction is needed for the upper bound.

For any progression and $1\leq U\leq V$,
\begin{equation}\label{eq:progression}
 \sum_{\substack{U\leq M\leq V\\M\equiv z\pmod q}}\frac1M
 \leq\frac1U+\frac1q\log(V/U).
\end{equation}
To prove it, keep the first progression term separately; for each
subsequent term at $t$, use $1/t\leq q^{-1}\int_{t-q}^{t}du/u$.
The integration intervals are disjoint and lie inside $[U,V]$.
If there are no progression terms, the bound is immediate.

If $2^A\leq x^{1/2}$, then
\[
 \frac q{M_0}=\frac{2^{b+1}3^n}{x}\leq2x^{-17/50}\leq1
\]
for large $x$.  Here $b\leq A$ and $3^n\leq x^{4/25}$, since
$n\leq L/(10\log2)$ and $3^5<2^8$.
Because $M_1/M_0=2^b$ and $b\geq1$, \eqref{eq:progression} gives
\begin{equation}\label{eq:low}
 c_{n,\mathbf a}(z)\leq C_2 b2^{-A}.
\end{equation}

If $2^A>x^{1/2}$, set $V_*=3^{m-1}2^{-(A-b)}M$.
The previous crossing inequality gives $V_*>x-F_{m-1}$, and therefore
$F_{m-1}=o(V_*)$ uniformly.  The last valuation is
\[
 b=\vnu(3(V_*+F_{m-1})+1).
\]
The positive integer inside the valuation is divisible by $2^b$, so
$2^b\leq3(V_*+F_{m-1})+1\leq5V_*$ for large $x$.
Consequently $M\geq(3/5)3^{-m}2^A$.  Apply
\eqref{eq:progression} with lower endpoint
$\max(M_0,(3/5)3^{-m}2^A)$ and upper endpoint $M_1$.
If they are out of order the contribution is zero; otherwise their ratio
is at most $2^b$.  We obtain
\begin{equation}\label{eq:high}
 c_{n,\mathbf a}(z)\leq C_3(3^n2^{-A}+b2^{-A}).
\end{equation}

The weights must now be summed before weakening them.  Exactly,
\[
 \sum_{a_1,\ldots,a_m\geq1}b2^{-A}
 =\left(\sum_{a\geq1}2^{-a}\right)^{m-1}
    \sum_{b\geq1}b2^{-b}=2.
\]
For the additional term in \eqref{eq:high},
\begin{align*}
 3^n\sum_{2^A>x^{1/2}}2^{-A}
 &\leq3^nx^{-1/4}\sum_{a_1,\ldots,a_m\geq1}2^{-A/2}\\
 &=3^nx^{-1/4}(1+\sqrt2)^m
 \leq x^{4/25-1/4+1/100000}
 =x^{-8999/100000}.
\end{align*}
The last step uses $1+\sqrt2<e$ and $m\leq L/100000$.
All summands are nonnegative, so enlarging the actual finite tuple range
to all positive tuples is valid.  Summing \eqref{eq:low} and
\eqref{eq:high} proves $c_n(z)\leq C_4(2+x^{-8999/100000})$, uniformly
in every stated variable.
\end{proof}

The last displayed majorant in the printed proof of Lemma~5.3 is instead
$\sum_{a_1,\ldots,a_m\geq1}2^{-A/2}=(1+\sqrt2)^m$.
It is finite for a fixed $m$ but is not uniformly bounded when
$m=\lfloor L/100000\rfloor$ grows.  Keeping the sharper weights already
available in the preceding estimates proves exactly the uniform result
needed by that lemma.

To apply the uniform bound, first observe that the source range gives
$r=n-m\geq m$: the lower endpoint gives
$n\geq(\alpha-1)L/\beta+L^{4/5}$, whereas
$2m\leq(\alpha-1)L/50$ and $\beta<1$.
If $p_r$ is the law of Tao's
$r$-Syracuse random variable on $\Z/3^r\Z$, define the uniform lift of
its reduction to $\Z/3^m\Z$ by
\[
 (L_{m\to r}p_m)(z)=3^{m-r}p_m(z\bmod3^m)\quad(0\leq m\leq r).
\]
Reduction has fibres of size $3^{r-m}$, so this lift preserves mass.
Tao's projective identity says that reducing $p_r$ gives $p_m$.
His fine-scale mixing estimate controls
$\normone{p_r-L_{m\to r}p_m}$.  Proposition~\ref{prop:cn} now gives
the precise dual estimate
\[
 \left|\sum_{z\bmod3^r}c_n(z)
       \bigl(p_r(z)-(L_{m\to r}p_m)(z)\bigr)\right|
 \leq C\normone{p_r-L_{m\to r}p_m}.
\]
The coefficient therefore transfers the fine-scale mixing error with
a constant independent of $x$, $n$, and $E$.

\section{Reserve bounds and threshold reductions}\label{sec:versions}

Two further parts of the first-passage argument depend on retaining the
appropriate finite-scale information: lifting a congruence between offset
integers requires a size bound, and passing to logarithmic density requires
the finite dyadic weights and the correct range of a threshold function.

\subsection{A reserve for \texorpdfstring{$3$}{3}-adic separation}

The $3$-adic separation argument in \cite[Section~6]{Tao7} retains a
reserve of $0.99C^2\log n$. The following estimate gives an explicit
sufficient margin for this size argument. Put
\[
 a=\log(4/3),\qquad b=\log2,\qquad \rho_*=\frac{b}{4a}.
\]
For positive integers $a_1,\ldots,a_k$, write $S_j=a_1+\cdots+a_j$
and $l=S_k$.  Clearing the dyadic denominator of the reversed offset
produces the positive integer
\[
 U=\sum_{j=1}^k3^{j-1}2^{l-S_j}.
\]
Bounding $U$ below $3^n$ is what permits an equality modulo $3^n$
between two such integers to be lifted to an equality of integers.

\begin{proposition}[An explicit reserve bound]\label{prop:reserve}
Suppose $n\geq2$, $C>0$, and $\rho>\rho_*$ satisfy
\[
 S_j\geq2j-2-C(\sqrt{j\log n}+\log n)\quad(1\leq j\leq k),
 \qquad
 l\leq n\frac{\log3}{b}-\rho C^2\log n.
\]
Put $\gamma=\rho b-b^2/(4a)>0$.  Then
\begin{equation}\label{eq:reserve}
 \frac{U}{3^n}\leq\frac83 n^{1-\gamma C^2+bC}.
\end{equation}
In particular $U<3^n$ if $\gamma C^2-bC\geq3$.
\end{proposition}
\begin{proof}
For each $j$, the hypotheses give
\[
 \frac{3^{j-1}2^{l-S_j}}{3^n}
 \leq\frac43\exp\left(-\rho b C^2\log n+bC\log n
                  -aj+bC\sqrt{j\log n}\right).
\]
Complete the square:
\[
 -aj+bC\sqrt{j\log n}
 =-a\left(\sqrt j-\frac{bC\sqrt{\log n}}{2a}\right)^2
  +\frac{b^2C^2\log n}{4a}.
\]
Every summand is thus at most $\frac43 n^{-\gamma C^2+bC}$.
Since $k\leq l<n\log3/\log2<2n$, summation gives
\eqref{eq:reserve}.  Under the stated last condition its right-hand
side is at most $8/(3n^2)\leq2/3<1$.
\end{proof}

Numerically $\rho_*=0.602355\ldots$, so the coefficient $0.99$ used by
Tao satisfies the strict inequality required here. The displayed bound
quantifies the sufficient reserve; it does not assert that this threshold
is necessary for other estimates of the same sum.

Here is the connection with the stopping event used there.  If
$T=n\log3/\log2-C^2\log n$ and the accumulated valuation first exceeds
$T$ at index $k+1$, the source concentration condition for the last two
coordinates gives
\[
 S_{k+1}\leq T+1+C(\sqrt{\log n}+\log n)
 \leq T+6C\log n.
\]
For $C\geq600$ this is at most
$n\log3/\log2-0.99C^2\log n$, exactly the retained reserve.
In deriving the first inequality, use $S_k\leq T$ and
$a_k\geq1$ in
$|a_k+a_{k+1}-2|\leq C(\sqrt{\log n}+\log n)$.
The source has already restricted the index to
$k=n\log3/(2\log2)+O(C\sqrt{n\log n})$, so $k\geq1$ for sufficiently
large $n$ at fixed $C$.  The prefix bound in
Proposition~\ref{prop:reserve} allows the additive two arising from
the source's interval indexing: for $j\geq2$, its interval starting
at index one gives $S_j\geq2(j-1)-C(\sqrt{(j-1)\log n}+\log n)$;
the case $j=1$ follows from $a_1\geq1$.

The role of the reserve is to ensure that congruent offset integers in
this size range are equal as integers.  Fourier decay additionally uses
the analytic estimates in Tao's argument.

\subsection{Dyadic weights and diverging thresholds}

The density factor in v7 Section~1.2 assigns mass $2^{-a}$ to
$\vnu(n)=a$.  Proposition~\ref{prop:dyadic} shows that the absolute
logarithmic density is $2^{-a-1}$; already $a=0$ distinguishes the
formulas.  The source prints the partial sum as $1-2^{-A_0}$;
the corrected finite partial sum is $1-2^{-A_0-1}$.
These two expressions tend to one, so this slip does not alter the limiting
conclusion, but it must be corrected in a finite-weight calculation.

At the end of Section~3, the source sets
$\widetilde f(X)=\inf_{m\geq X,\ m\text{ odd}}f(m)$ and applies it to a
sum over $m\leq X$.  This lower envelope does not control the earlier
values.  For a counterexample to the displayed bound, take $f(1)=0$ and
$f(m)=e^m$ for odd
$m\geq3$.  The source's bad harmonic sum is exactly $1$, since
$\Syr_{\min}(1)=1$ and $\Syr_{\min}(m)\leq m<e^m$ for all other $m$.
Its proposed bound is $O(\log X/X^c)$ and tends to zero.
The fixed-threshold argument of Corollary~\ref{cor:tao} supplies the
complete conclusion, including the strict inequality and arbitrary
real-valued nonmonotone functions, without that display.

\subsection{Source locations}

The following locators refer to \texttt{collatz.tex} in arXiv v7; the
source manifest records its SHA-256 identity.  Line numbers refer to that
TeX source, alongside the section and proposition numbers, rather than
to the journal pagination.

\begin{center}\small
\begin{tabular}{@{}>{\raggedright\arraybackslash}p{0.25\textwidth}>{\raggedright\arraybackslash}p{0.25\textwidth}>{\raggedright\arraybackslash}p{0.42\textwidth}@{}}
\toprule
Source locus & Estimate & Corresponding result\\
\midrule
Section 1.2, line 206 & Dyadic density factor & Exact finite transport,
Propositions~\ref{prop:dyadic} and \ref{prop:sampling}.\\[4pt]
Section 3, lines 591--593 & Tail envelope used on the wrong range &
Counterexample to the display; full replacement in
Corollary~\ref{cor:tao}.\\[4pt]
Proposition 5.2 proof, lines 759--762 & Additive input buffer &
Multiplicative buffer, Proposition~\ref{prop:buffer}.\\[4pt]
Lemma 5.3 proof, lines 868, 899, 913--914 & Previous-iterate exponent;
nonuniform final sum & Exact affine coefficient and uniformly summable
weights, Proposition~\ref{prop:cn}.\\[4pt]
Section 6, lines 989, 997, 1080--1084 & Retained offset reserve &
Explicit sufficient size bound,
Proposition~\ref{prop:reserve}.\\
\bottomrule
\end{tabular}
\end{center}


\section*{Computational materials and authorship}
The accompanying dependency-free Python certificate checks finite clock
identities, harmonic-weight transport, first-passage composition including
failure, and the algebraic constants in the estimates using exact
arithmetic.  Source identities and proof dependencies accompany the
certificate.  This work was developed through human--LLM collaboration
using ChatGPT 5.6 Sol, Ultra mode, in Codex, including source comparison,
proof development, and computational checks.
\begin{thebibliography}{9}
\bibitem{Allouche}
J.-P. Allouche, \emph{Sur la conjecture de ``Syracuse--Kakutani--Collatz''},
S\'eminaire de Th\'eorie des Nombres de Bordeaux (1978--1979),
expos\'e 9, 9-01--9-15, 1979.
\url{https://eudml.org/doc/182044}.

\bibitem{Everett}
C. J. Everett, \emph{Iteration of the number-theoretic function
$f(2n)=n$, $f(2n+1)=3n+2$}, Advances in Mathematics
\textbf{25} (1977), no.~1, 42--45.

\bibitem{Korec}
I. Korec, \emph{A density estimate for the $3x+1$ problem},
Mathematica Slovaca \textbf{44} (1994), no.~1, 85--89.
\url{https://dml.cz/handle/10338.dmlcz/133225}.

\bibitem{Tao7}
T. Tao, \emph{Almost all orbits of the Collatz map attain almost bounded
values}, arXiv:1909.03562v7, 16 July 2026. First submitted 8 September 2019.
\url{https://arxiv.org/abs/1909.03562v7}.

\bibitem{TaoJournal}
T. Tao, \emph{Almost all orbits of the Collatz map attain almost bounded
values}, Forum of Mathematics, Pi \textbf{10} (2022), e12, 56 pp.

\bibitem{Terras}
R. Terras, \emph{A stopping time problem on the positive integers},
Acta Arithmetica \textbf{30} (1976), no.~3, 241--252.
\url{https://doi.org/10.4064/aa-30-3-241-252}.
\end{thebibliography}

\end{document}
